The problems documented in Section~\ref{sec:description} point to a broad statistical fragility. Here, we identify the common structure underlying these failures: the non-identifiability of interpretability queries.

Behavior-based approaches that study input--output relationships (e.g., feature attributions, behavioral testing) are fundamentally limited by \textbf{underspecification}: multiple, distinct explanations can equally well account for the same input--output patterns \citep{jacovi2021contrastive,rogers2021primer,hagendorff2023machine}. Similar observations in cognitive science motivated the development of brain imaging as a complement to purely behavioral data, with the goal of measuring neural computation and thereby obtaining objective, measurable, and more generalizable quantities \citep{kosslyn1999if,logothetis2008we,churchland1988perspectives}. AI interpretability has followed a related trajectory moving toward analyzing internal computation \citep{mueller2024questrightmediatorhistory}. However, predictive approaches based on internal states (probing, SAEs) inherit standard machine-learning pathologies such as overfitting and poor generalization \citep{belinkov-2022-probing}. These failure modes are instances of \textbf{underspecification}: many predictive models can fit the training data equally well, leaving it unclear which ones posit generalizable causal mechanisms \citep{exp_unders,10.5555/3586589.3586815}.

Causal approaches, introduced in response to the shortcomings of predictive methods, appear at first to provide the scientific rigor needed for generalizable explanations. However, AI systems are large, distributed systems with many interacting components, which gives rise to redundant and context-dependent causal pathways~\citep{FrankleC19}. This creates \textbf{overdetermination}, where multiple distinct causal mechanisms are each independently sufficient to produce the same behavior \citep{doi:10.1073/pnas.91.11.5033,loosemore2012complex,sarkar2022explainableairacemodel}. Then, finding mechanistic stories within complex computational systems can become \textit{too easy}: many different, incompatible explanations can be produced for the same phenomenon \citep{lindsay2023testing,melouxeverything}.


\medskip
\noindent\textbf{Identifiability.}
These failure modes can be formalized using the concept of identifiability. Informally, identifiability is the property of a statistical inference task stating that the parameters (explanatory variables) of a statistical model can be uniquely recovered from available observations~\citep{casella2024statistical}. Identifiability is typically a prerequisite for reliable inference in the natural sciences; without it, inferred explanations remain ambiguous. Therefore, substantial work in statistics, unsupervised learning, and causal inference has focused on characterizing identifiability conditions and designing identifiable tasks~\citep{casella2024statistical,Allman2009,Locatello2019,Khemakhem2020,Shpitser2008}.

For interpretability, both underspecification and overdetermination produce non-identifiability, explaining most of the statistical fragilities:
(i) \textbf{Poor generalization:} when multiple explanations fit the observed data equally well, their explanatory claims can diverge arbitrarily on unseen data. Selecting among these explanations, therefore, depends on arbitrary inductive biases that are rarely validated.
(ii) \textbf{Sensitivity to design choices:} non-identifiability implies a manifold of explanations that achieve a \textit{good fit}. Different algorithmic choices (datasets, optimization procedures, hyperparameters) traverse this manifold differently, and thus produce different explanations.
(iii) \textbf{False discovery:} when explanations are non-identifiable, the probability of recovering a spurious explanation that happens to fit the data increases with the size and complexity of the hypothesis space.

Currently, identifiability is just a conceptual analogy, because interpretability has not yet been formalized as an explicit statistical inference task. Making this formal connection and casting interpretability queries as well-specified statistical estimation problems is a necessary first step toward developing methods whose limitations and assumptions can be explicitly characterized.